\subsubsection*{Ray characters and the two Frobenius prescriptions}
\label{subsec:aux-coordinates}
\hypertarget{source-auxiliary-coordinates}{}

The variables of the group-ring determinant must be realized by ray characters. A new prime must satisfy two Frobenius prescriptions, and both prescriptions have to survive the restrictions imposed at the distinguished twist prime.

Fix the auxiliary CM field \(E\), its conjugation \(c\) (with the chosen
extension of order two), the coefficient prime \(\ell\), and
\(L=\ell^M\), \(M\ge1\). Write \(T\) for the places over the
distinguished twist prime. Its residue characteristic is prime
to \(2L\). We first construct an open normal subgroup \(H\)
of \(G_E\), containing inertia above \(T\), and a finite set
\(S\) such that each split place \(P\notin S\) whose
decomposition group lies in \(H\) admits the coordinate
character of Section~\ref{subsec:retained-twist}.

\begin{proposition}[Prescribed Frobenius elements for the ray coordinates]\label{prop:ray-coordinate-pairs}
For the prescribed finite list of Frobenius conjugacy
conditions, the two groups of split auxiliary pairs can
be chosen outside any enlarged finite avoidance set.
Their anticyclotomic coordinates are surjective on the local
inertia of their own pair, unramified off that pair, and
trivial at decomposition groups above \(\ell\). Writing \(\eta_i\) for the additive first-group coordinates and \(P_j\) for the marked primes of the second group, one has
\[
 \eta_i(\operatorname{Frob}_{P_j})=\delta_{ij}\quad\text{in }\mathbb Z/L.
\]
\end{proposition}

\begin{proof}
\proofstep{Kill the global unit obstruction before prescribing a ray coordinate.}
Choose ideal representatives for the finite class group of \(E\)
prime to \(T\). This is possible by weak approximation at \(T\)
applied to the valuations of each representative. Put their prime
supports and the primes above \(\ell\) in \(S\), disjoint from \(T\).
Every fractional ideal has a generator away from \(S\). Choose
generators \(u_1,\ldots,u_a\) for the \(S\)-units, including
generators for their torsion. Take for \(H\) the subgroup fixing
\(E(\mu_L,u_1^{1/L},\ldots,u_a^{1/L})\). This is a finite
Galois extension. At \(v\in T\) every radicand is a unit
and \(L\) is invertible. The roots of each defining polynomial
lift its simple roots over the maximal unramified extension, so
inertia at \(v\) lies in \(H\).

For \(P\) as above, \(L\mid (\mathrm NP-1)\) and every \(u_i\)
reduces to an \(L\)-th power. Choose a surjection
\[
 \omega:k(P)^\times\to\mathbb Z/L
\]
from the cyclic residue
group. For an ideal \(J\) prime to \(P\), take a generator
\(\alpha\) away from \(S\). Its order at \(P\) is zero, so
\[
 J\mapsto\omega(\bar\alpha)
\]
is defined. Two generators
differ by an \(S\)-unit, and multiplication of generators proves
additivity. A principal ideal generated by \(1\bmod P\) has
value zero. Ray class reciprocity now gives a continuous
character \(\xi_P:G_E\to\mathbb Z/L\), unramified away from
\(P\). Its image on inertia at \(P\) is the image of \(k(P)^\times\):
on principal ideals prime to \(P\) the character evaluates the
residue character, and every residue unit lifts to an element
prime to \(P\).

At every \(v\mid\ell\) take \(1\) for a
generator of \(v\) away from \(S\). Both Frobenius and inertia
therefore act trivially there. Set
\[
       d_P(g)=\xi_P(g)-\xi_P(cgc^{-1}).
\]
Since \(P\ne cP\), only one summand ramifies at each member of
the pair. For the resulting difference character \(d_P\), the two inertia images are
\[
 d_P(I_P)=d_P(I_{cP})=\mathbb Z/L,
\]
including at \(\ell=2\).
This gives the desired anticyclotomic character, trivial on
decomposition groups above \(\ell\). \cite{milne-cft}

The universal coordinates in the theta data can be fixed once
per unordered pair. For each pair take a character with these
four properties (anti-invariance, triviality above \(\ell\),
ramification confined to the pair and surjectivity on its two
inertia groups) if one exists, and the trivial character otherwise.
Use this choice for the coordinate at \(P\) in every
conjugation-stable set containing it. Thus the ray construction
certifies existence at every admitted pair and the actual universal
coordinate there has the required properties. No uniqueness of a
ray character is required.
Restrict the chosen character to the local level of each selected series.

Here is the form of prime selection needed for the Galois
determinant. Let \(U\) be any open normal subgroup and let \(z\)
be the prescribed tame element at \(T\). Let
\[
       W=U\cap H\cap G_{E(\mu_{2L})}.
\]
We can prescribe the Frobenius of each first-half prime \(P_b\)
to be conjugate to \(z\) modulo \(W\), and the Frobenius of
each second-half prime \(P'_a\) to be conjugate modulo \(U\)
to any specified \(h_a\in H\cap G_{E(\mu_{2L})}\). In addition
we can require
\begin{equation}\label{eq:aux-coordinate}
       d_{P_b}(\mathrm{Frob}_{P'_a})=\mathbf1_{a=b}
                         \quad\text{in }\mathbb Z/L .
\end{equation}

To prove this, apply Chebotarev over \(E\) in the finite
quotient by \(W\), avoiding any given finite set, the ramification
loci, their conjugates and the primes above \(2L\). Degree-one
primes suffice: in the prime count over a number field primes of
higher absolute residue degree up to norm \(x\) contribute
\(O_E(\sqrt{x}\log x)\). A degree-one prime unramified over
\(\mathbb Q\) splits over its restriction to the maximal real
subfield because an inert quadratic place would have even residue
degree. The element \(z\) lies in \(H\) and fixes \(\mu_{2L}\)
by the inertia condition on \(T\). The primes selected have
\(\mathrm NP_b\equiv1\bmod 2L\) and are admitted. Their
inertia groups lie in \(W\). After each choice exclude the
whole pair.

\proofstep{Multiply each later Frobenius by inertia before applying Chebotarev.}
The new Frobenius must realize the prescribed values on the earlier ray coordinates as well.

The map
\[
 W\to(\mathbb Z/L)^R
\]
given by the \(R\) first-half
characters is surjective. Inertia at \(P_b\) realizes every
multiple of the \(b\)-th basis vector, and the target is abelian.
Choose \(e_a\in W\) so that the coordinates of \(h_a e_a\)
are the \(a\)-th basis vector. Apply Chebotarev again, now
modulo
\[
 W\cap\bigcap_b\ker d_{P_b},
\]
to \(h_a e_a\).
These primes are admitted by the same norm and splitting tests,
satisfy the prescribed conjugacies modulo \(U\), and satisfy
\eqref{eq:aux-coordinate} for every Frobenius lift. All exclusions
can be enlarged before either choice. This proves the prime and
character part of the construction at every depth. Bounds on
the norms of the auxiliary primes do not enter the length
estimate. The numbers of pairs and the degree tests uniform in
\(M\) come from the residual Selmer argument below.
\end{proof}

